\introduction % Структурный элемент: ВВЕДЕНИЕ

В современном гражданском обороте Российской Федерации миллионы сделок ежедневно заключаются между физическими и юридическими лицами, причём существенная доля таких сделок относится к типовым массовым правоотношениям~--- розничной купле-продаже, бытовому подряду, возмездному оказанию услуг, найму жилых помещений, прокату, дарению, займу, сделкам с недвижимым имуществом. Согласно статье~432 Гражданского кодекса Российской Федерации, договор считается заключённым, если между сторонами достигнуто соглашение по всем его существенным условиям в требуемой форме; нарушение этих требований влечёт за собой признание договора незаключённым и сопряжённые с этим финансовые и временные потери для сторон. Между тем большинство участников гражданского оборота не имеет юридического образования, а квалифицированная юридическая помощь при подготовке бытовых договоров экономически малодоступна. Сложившиеся к настоящему моменту шаблонные конструкторы договоров не дают правовой оценки описанной сделки с опорой на действующее законодательство, а лишь подставляют пользовательские значения в заранее заготовленный текст из закрытого каталога сценариев; универсальные большие языковые модели работают со свободной формулировкой запроса, но систематически продуцируют правовые галлюцинации и не привязаны к актуальной редакции законодательства; существующие в правовой сфере системы генерации, дополненной поиском по внешней памяти (retrieval-augmented generation, RAG)~\cite{lewis2020rag}, в основном ограничиваются сценарием <<вопрос-ответ>> и не подготавливают полноценный проект договора. Совмещение массовой потребности в первичной правовой поддержке и сложившейся технологической базы~--- больших языковых моделей (large language models, LLM), методов RAG, открытых векторных баз данных, средств локального исполнения языковых моделей,~--- определяет актуальность настоящей работы.

Объект исследования~--- процесс автоматизированной подготовки проектов договоров по российскому гражданскому праву по словесному описанию планируемой сделки. Предмет исследования~--- архитектура и программная реализация интеллектуального агента, обрабатывающего произвольную пользовательскую формулировку сделки с опорой на проиндексированный корпус российского законодательства и с автоматическим контролем качества подготавливаемого проекта договора.

Цель работы~--- разработать интеллектуальный программный агент, способный по произвольному словесному описанию сделки определять тип соответствующего ей гражданского-правового договора, в диалоговом режиме собирать у пользователя недостающие для подготовки документа сведения и формировать проект договора в формате DOCX вместе с юридическими рекомендациями со ссылками на конкретные статьи нормативных правовых актов.

Для достижения поставленной цели в настоящей работе решаются следующие задачи:
\begin{enumerate}
    \item сформировать корпус нормативных правовых актов российского договорного права и реализовать его предобработку, разбиение на чанки и индексацию в векторной базе данных;
    \item спроектировать формальную модель агента в виде графа состояний с условными переходами, в узлах которого выполняются отдельные шаги обработки запроса (классификация типа сделки, проверка общих норм, поиск релевантных специальных норм, проверка достаточности данных, генерация проекта договора, его валидация и формирование рекомендаций);
    \item реализовать механизм диалогового уточнения недостающих сведений: при нехватке данных для согласования существенных условий договора агент формулирует пользователю целевой уточняющий вопрос и продолжает обработку после получения ответа;
    \item реализовать автоматическую итеративную проверку сгенерированного проекта договора на полноту существенных условий и непротиворечивость применимым нормам с возможностью повторной генерации при обнаружении замечаний;
    \item разработать модуль генерации файла договора в формате DOCX с надлежащим оформлением документа делового стиля;
    \item разработать веб-интерфейс пользователя и панель администратора для управления корпусом нормативных правовых актов и каталогом поддерживаемых типов сделок;
    \item обеспечить контейнерное развёртывание программной системы и её работоспособность на локальном оборудовании без передачи пользовательских данных во внешние сервисы;
    \item провести внешнюю экспертную оценку качества подготовленных системой артефактов с привлечением практикующих юристов.
\end{enumerate}

В работе применяются методы анализа предметной области (нормативно-правовой анализ положений Гражданского кодекса Российской Федерации, регламентирующих заключение договора; обзор существующих программных решений в области правовых технологий и научной литературы по применению больших языковых моделей к юридическим задачам); методы проектирования программных систем (формализация поведения агента как ориентированного графа состояний с условными переходами; декомпозиция функций системы на независимые узлы графа); методы разработки систем поиска, дополненной нейросетевыми эмбеддингами и реранкерами (RAG-архитектура с двухступенчатой выдачей и LLM-фильтром); методы эмпирической оценки качества (количественные метрики качества поиска на эталонном датасете, экспертная апробация конечных артефактов с использованием пятибалльной шкалы Лайкерта~\cite{likert1932}).

Научная новизна работы заключается в построении интеллектуального агента, объединяющего в одной программной системе три свойства: работу с произвольной формулировкой сделки от пользователя, опору на проиндексированный корпус российского законодательства и автоматический контроль качества подготавливаемого проекта договора. Подобное сочетание свойств на момент выполнения настоящей работы в открытой литературе и в практических программных продуктах в области правовых технологий не зафиксировано.

Практическая значимость работы определяется реализацией предложенного подхода в виде целостного программного продукта, применимого для подготовки договоров по бытовым сделкам гражданского оборота физическими лицами и субъектами малого предпринимательства. Внешняя экспертная апробация конечных артефактов с привлечением практикующих юристов подтверждает приемлемый уровень качества проектов договоров и сопровождающих рекомендаций.
